\documentclass[12pt]{article}

\usepackage{latexsym}
\usepackage[T1]{fontenc} 
\usepackage[utf8]{inputenc} 
\usepackage[spanish]{babel}
\usepackage{amsthm, amssymb, amsmath, amsfonts, bbm, mathtools} 
\usepackage{multirow, array}
\usepackage{stmaryrd}
\usepackage{xcolor}
\decimalpoint

\usepackage{wrapfig}
\usepackage{subcaption}

% no tocar nada de aqui
\providecommand{\abs}[1]{\lvert#1\rvert}
\usepackage[hidelinks]{hyperref}
\usepackage{enumerate} 
\usepackage{graphicx}

% entornos personalizados
\newtheorem{example}{Ejemplo}
\newtheorem{notacion}{Notación}
\newtheorem{defi}{Definición}
\newtheorem{obse}{Observación}
\newtheorem{lem}{Lema}

\newcommand{\sol}{\textit{Solución. }}
\renewcommand{\qed}{\hfill $\blacksquare$}
\newcommand{\eej}{\hfill $\square$}

%--------------------- Comandos personalizados ---------------------

% Simbolos de los Numeros
\newcommand{\D}{\mathfrak{D}}
\newcommand{\A}{\mathcal{A}}
\newcommand{\C}{\mathbb{C}}
\newcommand{\F}{\mathbb{F}}
\newcommand{\I}{\mathbb{I}}
\newcommand{\N}{\mathbb{N}}
\newcommand{\Q}{\mathbb{Q}}
\newcommand{\R}{\mathbb{R}}
\newcommand{\Z}{\mathbb{Z}}

% Signos de agrupacion
\newcommand{\pp}[1]{\left( #1 \right)}
\newcommand{\cc}[1]{\left[ #1 \right]}
\newcommand{\set}[1]{\left\{#1\right\}}

\newcommand{\ol}[1]{\overline{#1}}
\newcommand{\ceil}[1]{\left\lceil#1\right\rceil}
\newcommand{\floor}[1]{\left\lfloor#1\right\rfloor}
\newcommand{\ind}[1]{\mathbbm{1}_{\{#1\}}}
\newcommand{\norm}[1]{\left\Vert #1 \right\Vert}

\newcommand{\Text}[1]{\quad\text{#1}\quad}

\newcommand{\bs}{\backslash}
\renewcommand{\c}{\complement}
\newcommand{\beq}{bilipschitz }

% Flechas
\newcommand{\imp}{\Longrightarrow}
\newcommand{\Imp}{\quad\Longrightarrow\quad}
\newcommand{\ssi}{\Longleftrightarrow}
\newcommand{\Ssi}{\quad\Longleftrightarrow\quad}


\newcommand{\e}{\varepsilon}
\renewcommand{\d}{\delta}
\renewcommand{\O}{\Omega}
\renewcommand{\l}{\lambda}
\renewcommand{\a}{\alpha}
\renewcommand{\b}{\beta}
\renewcommand{\abs}[1]{\left\lvert#1\right\rvert}
%\renewcommand{\qedsymbol}{$\blacksquare$}


\newcommand{\Co}{\mathfrak C}
\newcommand{\T}{\mathbb T}
%\newcommand{\uno}{\mathbb 1}
%\newcommand{\ind}[1]{\mathbbm{1}_{\{#1\}}}
\newcommand{\uno}{\mathbbm{1}}


\renewcommand{\o}{\omega}

\newcommand{\x}{\mathrm x}
\newcommand{\y}{\mathrm y}
\newcommand{\Int}{\mathrm{int}}
\newcommand{\Ext}{\mathrm{ext}}
\newcommand{\diam}{\mathrm{diam}}
\newcommand{\cl}{\mathrm{cl}}
\newcommand{\angles}[2]{\left\langle #1, #2\right\rangle}

\textheight=25cm \textwidth=18cm \topmargin=-2cm \oddsidemargin=-1cm

\setlength{\parskip}{0.8em}  % Espaciado párrafos
\setlength{\parindent}{0pt}  % Sangría

\newcommand{\To}{\to}



\begin{document}
    \pagestyle{empty}
    \begin{center}
        {\textsc{Ecuaciones Diferenciales Parciales II} \\
        \medskip 
        \large{Tarea 1}}
    \end{center}
    
    \bigskip
    
    \begin{enumerate}
        \item Resolver las siguientes PDE's relacionadas a la ecuación de calor.

        \begin{enumerate}
            \item \[
            \begin{cases}
                2u_{xx} &= u_t,\quad0<x<4,~t>0,\\
                u(0,t) &= 0,\\
                u(4,t) &= 0,\\
                u(x,0) &= \begin{cases}
                    x,\quad   &0<x<2,\\
                    4-x,\quad &2<x<4.
                \end{cases}
            \end{cases}
            \]

            \sol Por el método de separación de variables supongamos que la solución $u$ de la ecuación se puedo factorizar de la forma $u(x,t)=v(x) w(t)$ con $v,w\neq0$. Entonces, tenemos que
            \[
            \begin{cases}
                u(0,t) =0 \Ssi v(0) w(t) =0 \Ssi  v(0)=0,\\
                u(4,t) =0 \Ssi v(4) w(t) =0 \Ssi v(4) = 0.
            \end{cases}
            \]
            Por otro lado, como $u(x,t)=v(x) w(t)$, entonces
            \[ 2u_{xx} = u_t \Ssi 2v_{xx}w = vw_t\Ssi 2\frac{v''}{v} = \frac{w'}{w}=\l, \]
            donde $\l$ es cualquier constante, ya que las variables $x$ y $t$ no tienen relación entre ellas. De lo anterior obtenemos
            \[
            \begin{cases}
                2v''-\l v &=0, \quad v(0)=v(4)=0,\\
                w'-\l w &=0.
            \end{cases}
            \]
            Supongamos que la solución de la primera ecuación es de la forma $v(x)=e^{rx}$ para alguna $r$. Entonces
            \begin{align*}
                2v''-\l v =0 &\Ssi2r^2e^{rx}-\l e^{rx}=0\\
                & \Ssi e^{rx}(2r^2-\l)=0\\
                &\Ssi 2r^2-\l=0\\
                &\Ssi r^2=\frac{\l}{2}.
            \end{align*}
            Si $\l>0$, entonces la solución general es
            \[ v(x) = A e^{\sqrt{\tfrac{\l}{2}}x} + B e^{-\sqrt{\tfrac{\l}{2}}x}, \]
            con $A, B\in\R$ constantes cualesquiera. Sin embargo, para que se cumplan las condiciones de frontera se debe cumplir que
            \[
            \begin{cases}
                v(0)=A e^{\sqrt{\tfrac{\l}{2}}\cdot 0} + B e^{-\sqrt{\tfrac{\l}{2}}\cdot 0} =0,\\
                v(4)=A e^{\sqrt{\tfrac{\l}{2}}\cdot 4} + B e^{-\sqrt{\tfrac{\l}{2}}\cdot4}=0,
            \end{cases}
            \]
            que es equivalente al sistema
            \[
            \begin{cases}
                A + B  =0,\\
                A e^{\sqrt{8\l}} + B e^{-\sqrt{8\l}}=0.
            \end{cases}
            \]
            Despejando en la primera ecuación tenemos que $B=-A$ y sustituyendo en la segunda tenemos
            \[ A e^{\sqrt{8\l}} -A e^{-\sqrt{8\l}}=A\pp{e^{\sqrt{8\l}} -e^{-\sqrt{8\l}}}=0. \]
            Como pedimos que $v\neq 0$, entonces $A\neq 0$ y consideramos los $\l$ tales que $e^{\sqrt{8\l}} =e^{-\sqrt{8\l}}$, pero esto sólo ocurre para $\l=0$, lo cual es una contradicción a que $\l>0$. Por lo tanto, 
            \[ 2v''-\l v =0, \quad v(0)=v(4)=0, \]
            no tiene solución no trivial para $\l>0$.

            Para el caso $\l=0$ se tiene que $2v''-0\cdot v = 2v'' =0$, entonces $v(x)=Ax+B$, con $A,B\in\R$ constantes cualesquiera. Para que se cumpla que $v(0)=v(4)=0$, se debe cumplir que $A\cdot0+B=B=0$ y $v(4)=A\cdot 4+B=0$, esto es, $A=B=0$ que implica que $v=0$, pero esta solución no nos interesa.

            Por otro lado, para el caso $\l<0$, se tiene que $v(x)=e^{rx}$ donde $r^2=\l/2$. Podemos escribir $\l=-\a^2$ con $\a>0$, entonces 
            \[r^2=\l/2=-\a^2/2\Ssi r=\pm i\frac{\a}{\sqrt{2}}.\]
            Se sigue que la solución general es
            \[v(x) = A\cos\pp{\frac{\a}{\sqrt{2}} x} + B\sin\pp{\frac{\a}{\sqrt{2}} x},\]
            pero dado que $v(0)=v(4)=0$, entonces se debe cumplir que
            \[
            \begin{cases}
                A\cos\pp{\frac{\a}{\sqrt{2}}\cdot 0} + B\sin\pp{\frac{\a}{\sqrt{2}}\cdot 0} = 0,\\
                A\cos\pp{\frac{\a}{\sqrt{2}}\cdot 4} + B\sin\pp{\frac{\a}{\sqrt{2}}\cdot 4} =0,
            \end{cases}
            \Ssi
            \begin{cases}
                A = 0,\\
                B\sin\pp{2\sqrt{2}\cdot\a} =0.
            \end{cases}
            \]
            Como $A=0$ y queremos que $v\neq 0$, entonces debemos considerar los $\a>0$ tales que $2\sqrt{2}\cdot \a$ anulen la función seno, estos son todos los $\a_k$ con $k\in\N$ tales que
            \begin{align*}
                2\sqrt{2}\cdot \a_k = k\pi,\Ssi \a_k=\frac{k\pi\sqrt{2}}{4}.
            \end{align*}
            Por lo tanto,
            \[ v_k(x) = \sin\pp{\frac{\a_k}{\sqrt{2}}\cdot x}= \sin\pp{\frac{k\pi x}{4}}, \]
            son soluciones para todo $k\in\N$.

            Por otro lado, resolviendo $w'-\l_k w=0$ donde $\l_k=-{\a_k}^2$ tenemos
            \begin{align*}
                w'-\l_k w=0 &\Imp \frac{w'}{w}=-{\a_k}^2\\
                &\Imp \frac{w'}{w}=-\pp{\frac{k\pi\sqrt{2}}{4}}^2\\
                &\Imp w_k(t)= A_ke^{-\frac{(k\pi)^2}{8}t}.
            \end{align*}
            Entonces la solución de la ecuación original es
            \[ u(x,t) = \sum_{k\in\N}v_k(x)w_k(t)=\sum_{n\in\N}C_k\sin\pp{\frac{k\pi x}{4}}e^{-\frac{(k\pi)^2}{8}t}, \]
            donde las $C_k$ son contantes. Sin embargo, se debe cumplir que
            \[u(x,0) = f(x):=\begin{cases}
                x,\quad   &0<x<2,\\
                4-x,\quad &2<x<4,
            \end{cases} \]
            esto es,
            \[
            \sum_{n\in\N}C_k\sin\pp{\frac{k\pi x}{4}}= f(x). \]
            Por la expansión en series de Fourier se sigue que los $C_k$ son los coeficientes de la serie de Fourier de senos de $f$ en $[0,4]$. Entonces
            \begin{align*}
                C_k &= \frac{2}{4}\int_0^4 f(x)\sin\pp{\frac{k\pi x}{4}}dx\\
                &= \frac{1}{2}\int_0^2 x\sin\pp{\frac{k\pi x}{4}}dx+\frac{1}{2}\int_2^4 (4-x)\sin\pp{\frac{k\pi x}{4}}dx\\
                &= \frac{1}{2}\int_0^2 x\sin\pp{\frac{k\pi x}{4}}dx+2\int_2^4\sin\pp{\frac{k\pi x}{4}}dx-\frac{1}{2}\int_2^4x\sin\pp{\frac{k\pi x}{4}}dx\\
                &= \frac{8}{(k\pi)^2}\int_0^{k\pi/2} \xi\sin(\xi)d\xi+\frac{8}{k\pi}\int_{k\pi/2}^{k\pi}\sin(\xi)d\xi-\frac{8}{(k\pi)^2}\int_{k\pi/2}^4\xi\sin(\xi)d\a\qquad (\xi=\tfrac{k\pi x}{4}).
            \end{align*}
            Luego, integrando por partes y usando las integrales de funciones trigonométricas conocidas tenemos que
            \begin{align*}
                C_k &= \frac{8}{(k\pi)^2}\int_0^{k\pi/2} \xi\sin(\xi)d\xi+\frac{8}{k\pi}\int_{k\pi/2}^{k\pi}\sin(\xi)d\xi-\frac{8}{(k\pi)^2}\int_{k\pi/2}^{k\pi}\xi\sin(\xi)d\a\\
                &= \frac{8}{(k\pi)^2}\int_0^{k\pi/2} \xi\sin(\xi)d\xi+\frac{8}{k\pi}\int_{k\pi/2}^{k\pi}\sin(\xi)d\xi-\frac{8}{(k\pi)^2}\int_{k\pi/2}^{k\pi}\xi\sin(\xi)d\a\\
                &= \frac{8}{(k\pi)^2}\cc{-{\xi}\cos\pp{\xi}+\sin\pp{\xi}}_{0}^{k\pi/2}-\frac{8}{k\pi}\cc{\cos\pp{\xi}}_{k\pi/2}^{k\pi}-\frac{8}{(k\pi)^2}\cc{-{\xi}\cos\pp{\xi}+\sin\pp{\xi}}_{k\pi/2}^{k\pi}\\
                &= \frac{8}{(k\pi)^2}\cc{-\frac{k\pi}{2}\cos\pp{\frac{k\pi}{2}}+\sin\pp{\frac{k\pi}{2}}}
                -\frac{8}{k\pi}\cc{(-1)^k-\cos\pp{\frac{k\pi}{2}}}
                \\&\qquad-\frac{8}{(k\pi)^2}\cc{-{k\pi}(-1)^{k}+\frac{k\pi}{2}\cos\pp{\frac{k\pi}{2}}-\sin\pp{\frac{k\pi}{2}}}\\
                &= \frac{8}{(k\pi)^2}\cc{-\frac{k\pi}{2}\cos\pp{\frac{k\pi}{2}}+\sin\pp{\frac{k\pi}{2}}+{k\pi}(-1)^{k}-\frac{k\pi}{2}\cos\pp{\frac{k\pi}{2}}+\sin\pp{\frac{k\pi}{2}}}
                \\&\qquad-\frac{8}{k\pi}\cc{(-1)^k-\cos\pp{\frac{k\pi}{2}}}\\
                &= \frac{8}{(k\pi)^2}\cc{-k\pi\cos\pp{\frac{k\pi}{2}}+2\sin\pp{\frac{k\pi}{2}}+{k\pi}(-1)^{k}}-\frac{8}{k\pi}\cc{(-1)^k-\cos\pp{\frac{k\pi}{2}}}.
            \end{align*}
            Si $k$ es par, entonces se puede escribir de la forma $k=2n$ con $n\in\N$, entonces
            \begin{align*}
                C_k &= \frac{8}{(2n\pi)^2}\cc{-2n\pi\cos\pp{n\pi}+2\sin\pp{n\pi}+{2n\pi}(-1)^{2n}}-\frac{8}{2n\pi}\cc{(-1)^{2n}-\cos\pp{n\pi}}\\
                &= \frac{2}{(n\pi)^2}\cc{-2n\pi(-1)^{n}+{2n\pi}}-\frac{4}{n\pi}\cc{1-(-1)^n}\\
                &= \frac{2}{(n\pi)^2}\cc{-2n\pi(-1)^{n}+{2n\pi}}-\frac{4}{n\pi}\cc{1-(-1)^n}\\
                &= -\frac{4}{n\pi}\cc{(-1)^{n}-1}+\frac{4}{n\pi}\cc{(-1)^n-1}\\
                &= 0.
            \end{align*}
            Si $k$ es impar, se puede escribir de la forma $k=2n-1$ con $n\in\N$, entonces
            \begin{align*}
                C_{2n-1} &= \frac{8}{[(2n-1)\pi]^2}\cc{-(2n-1)\pi\cos\cc{\frac{(2n-1)\pi}{2}}+2\sin\cc{\frac{(2n-1)\pi}{2}}-{(2n-1)\pi}}\\&\quad\quad+\frac{8}{k\pi}\cc{1+\cos\cc{\frac{(2n-1)\pi}{2}}}\\
                &= \frac{8}{[(2n-1)\pi]^2}\cc{2(-1)^{n+1}-{(2n-1)\pi}}+\frac{8}{(2n-1)\pi}\\
                &= -\frac{16(-1)^{n}}{[(2n-1)\pi]^2}.
            \end{align*}
            Por lo tanto
            \[ u(x,t) = -\sum_{n\in\N}(-1)^{n}\cc{\frac{4}{(2n-1)\pi}}^2\sin\cc{\frac{(2n-1)\pi x}{4}}e^{-\frac{[(2n-1)\pi]^2}{8}t}. \] \eej

            \item \[
            \begin{cases}
                \a^2u_{xx} &= u_t,\quad 0<x<L,~t>0,\\
                u(0,t) &= 0,\\
                u(L,t) &= 0,\\
                u(x,0) &= 2\sin\pp{\tfrac{5\pi x}{2L}}-4\sin\pp{\tfrac{9\pi x}{2L}},\quad 0<x<L.
            \end{cases}
            \]

            \sol Por el método de separación de variables, supongamos que la solución se descompone como $u(x,t)=v(x)w(t)$ con $v,w\neq 0$. Tenemos que
            \[
            \begin{cases}
                u(0,t) &= 0\\
                u(L,t) &= 0\\
            \end{cases}
            \Ssi
            \begin{cases}
                v(0)w(t) &= 0\\
                v(L)w(t) &= 0\\
            \end{cases}
            \Ssi
            \begin{cases}
                v(0) &= 0\\
                v(L) &= 0\\
            \end{cases}.
            \]
            Por otro lado
            \[ \a^2u_{xx} = u_t\Ssi \a^2 v''w = vw'\Ssi \a^2 \frac{v''}{v} = \frac{w'}{w}=\l,\]
            donde $\l$ es cualquier constante ya que las variables $x$ y $t$ son independientes entre ellas. Entonces tenemos el siguiente sistema de edos
            \[
            \begin{cases}
                \a^2v''-\l v=0,\quad v(0)=v(L)=0,\\
                w'-\l w=0.
            \end{cases}
            \]
            Tenemos que $v(x)=e^{rx}$ con $r$ alguna constante, entonces
            \[ \a^2v''-\l v=0\Ssi \a^2r^2e^{rx}-\l e^{rx}=0 \Ssi \a^2r^2-\l =0 \Ssi r^2 =\frac{\l}{\a^2}. \]
            Nuevamente, si $\l=0$, entonces la ecuación es $v''=0$ con solución general $v(x)=Ax+B$ pero para que se cumpla que $v(0)=v(L)=0$ se tiene que cumplir que $A=B=0$, entonces $v=0$, lo cual no queremos. Por otro lado, si $\l>0$, entonces $r=\pm\sqrt{\l}/\abs{a}$ y
            \[v(x)=Ae^{\tfrac{\l}{|\a|}x}+B e^{-\tfrac{\l}{|\a|}x}.\]
            Luego,
            \[v(0)=v(L)=0\Ssi
            \begin{cases}
                A+B =0,\\
                Ae^{\tfrac{\l}{|\a|}L}+B e^{-\tfrac{\l}{|\a|}L}=0,
            \end{cases}
            \Ssi
            \begin{cases}
                B=-A\\
                Ae^{\tfrac{\l}{|\a|}L}=A e^{-\tfrac{\l}{|\a|}L}
            \end{cases},
            \]
            lo anterior se cumple solamente cuando $A=B=0$ lo cual es equivalente a $v=0$ que no queremos que ocurra.

            Por otro lado, si $\l<0$, entonces $r=\pm i\sqrt{-\l}/|\a|$, se sigue que
            \begin{align*}
                v(x) &= Ae^{\frac{i\sqrt{-\l}}{|\a|}}+Be^{-\frac{i\sqrt{-\l}}{|\a|}}\\
                &=A\cos\pp{\frac{\sqrt{-\l}}{|\a|}x}+B\sin\pp{\frac{\sqrt{-\l}}{|\a|}x}.
            \end{align*}
            Luego,
            \begin{align*}
                v(0)=v(L)=0 &\Ssi \begin{cases}
                    A\cos\pp{\frac{\sqrt{-\l}}{|\a|}\cdot 0}+B\sin\pp{\frac{\sqrt{-\l}}{|\a|}\cdot 0}=0,\\
                    A\cos\pp{\frac{\sqrt{-\l}}{|\a|}L}+B\sin\pp{\frac{\sqrt{-\l}}{|\a|}L}=0.
                \end{cases}\\
                &\Ssi\begin{cases}
                    A=0,\\
                    B\sin\pp{\frac{\sqrt{-\l}}{|\a|}L}=0.
                \end{cases}
            \end{align*}
            Como queremos que $v\neq 0$, entonces
            \[ B\sin\pp{\frac{\sqrt{-\l}}{|\a|}L}=0,~ B\neq 0\Ssi \frac{\sqrt{-\l}}{|\a|}L=n\pi,~n\in\N \Ssi \l=-\pp{\frac{\a n\pi}{L}}^2. \]
            Sean $\l_n=-\pp{\a n\pi/L}^2$, entonces
            \[ v_n(x)=B_n\sin\pp{\frac{\sqrt{-\l_n}}{|\a|}x}=B_n\sin\pp{\frac{\sqrt{\pp{\a n\pi/L}^2}}{|\a|}x}=B_n\sin\pp{\frac{n\pi x}{L}} \]
            son soluciones.
            Para resolver ${w_n}'-\l_n w_n=0$, tenemos que
            \begin{align*}
                {w_n}'-\l_n w_n=0 &\Imp \frac{{w_n}'}{w_n} = \l_n,\\
                &\Imp w_n(t)= A_ne^{\l_nt} = A_ne^{-\pp{\frac{\a n\pi}{L}}^2t}.
            \end{align*}
            Por lo tanto,
            \[ u(x,t)=\sum_{n\in\N}v_n(x)w_n(t)=\sum_{n\in\N}C_n \sin\pp{\frac{n\pi x}{L}}e^{-\pp{\frac{\a n\pi}{L}}^2t}.\]
            Luego queremos que 
            \[ u(x,0)=\sum_{n\in\N}C_n \sin\pp{\frac{n\pi x}{L}}=2\sin\pp{\tfrac{5\pi x}{2L}}-4\sin\pp{\tfrac{9\pi x}{2L}}.\]
            Por la expansión de la serie de Fourier de senos tenemos que
            \begin{align*}
                C_n &= \frac{2}{L}\int_0^L\cc{2\sin\pp{\tfrac{5\pi x}{2L}}-4\sin\pp{\tfrac{9\pi x}{2L}}}\sin\pp{\tfrac{n\pi x}{L}}dx\\
                &= \frac{4}{L}\int_0^L\sin\pp{\tfrac{5\pi x}{2L}}\sin\pp{\tfrac{n\pi x}{L}}dx-\frac{8}{L}\int_0^L\sin\pp{\tfrac{9\pi x}{2L}}\sin\pp{\tfrac{n\pi x}{L}}dx.
            \end{align*}
            Usando que
            \[ \sin(Ax)\sin(Bx)=\frac{\cos[(A-B)x]-\cos[(A+B)x]}{2}, \]
            tenemos que
            \begin{align*}
                C_n &= \frac{2}{L}\int_0^L\cc{\cos\pp{\tfrac{5\pi-2n\pi}{2L}x}-\cos\pp{\tfrac{5\pi+2n\pi}{2L}x}}dx\\
                &\qquad-\frac{8}{L}\int_0^L\cc{\cos\pp{\tfrac{9\pi-2n\pi}{2L}x}-\cos\pp{\tfrac{9\pi+2n\pi}{2L}x}}dx\\
                &= \frac{2}{L}\int_0^L\cos\pp{\tfrac{5\pi-2n\pi}{2L}x}dx-\frac{2}{L}\int_0^L\cos\pp{\tfrac{5\pi+2n\pi}{2L}x}dx\\
                &\qquad-\frac{8}{L}\int_0^L\cos\pp{\tfrac{9\pi-2n\pi}{2L}x}dx+\frac{8}{L}\int_0^L\cos\pp{\tfrac{9\pi+2n\pi}{2L}x}dx.
            \end{align*}
            Sabemos que
            \[ \int\cos(Ax)dx=\tfrac{1}{A}\sin(Ax)+C. \]
            Entonces
            \begin{align*}
                C_n &= \frac{4}{\pi(5-2n)}\cc{\sin\pp{\tfrac{5\pi-2n\pi}{2L}x}}_0^L-\frac{4}{\pi(5+2n)}\cc{\sin\pp{\tfrac{5\pi+2n\pi}{2L}x}}_0^L\\
                &\qquad-\frac{16}{\pi(9-2n)}\cc{\sin\pp{\tfrac{9\pi-2n\pi}{2L}x}}_0^L+\frac{16}{\pi(9+2n)}\cc{\sin\pp{\tfrac{9\pi+2n\pi}{2L}x}}_0^L\\
                &= \frac{4}{\pi(5-2n)}\sin\pp{\tfrac{5-2n}{2}\pi}-\frac{4}{\pi(5+2n)}\sin\pp{\tfrac{5+2n}{2}\pi}\\
                &\qquad-\frac{16}{\pi(9-2n)}\sin\pp{\tfrac{9-2n}{2}\pi}+\frac{16}{\pi(9+2n)}\sin\pp{\tfrac{9+2n}{2}\pi}\\
                &= \frac{4}{(2n-5)\pi}\sin\pp{\tfrac{2n-5}{2}\pi}-\frac{4}{(2n+5)\pi}\sin\pp{\tfrac{2n+5}{2}\pi}\\
                &\qquad-\frac{16}{(2n-9)\pi}\sin\pp{\tfrac{2n-9}{2}\pi}+\frac{16}{(2n+9)\pi}\sin\pp{\tfrac{2n+9}{2}\pi}\\
                &= \frac{4}{(2n-5)\pi}(-1)^{n+1}+\frac{4}{(2n+5)\pi}(-1)^{n+1}-\frac{16}{(2n-9)\pi}(-1)^{n+1}-\frac{16}{(2n+9)\pi}(-1)^{n+1}\\
                &= \frac{4}{\pi}(-1)^{n+1}\pp{\frac{1}{2n-5}+\frac{1}{2n+5}-\frac{4}{2n-9}-\frac{4}{2n+9}}\\
                &= \frac{4}{\pi}(-1)^{n+1}\pp{\frac{4n}{4n^2-5^2}-4\frac{4n}{4n^2-9^2}}\\
                &= \frac{16n}{\pi}(-1)^{n+1}\pp{\frac{1}{4n^2-5^2}-\frac{4}{4n^2-9^2}}\\
                &= \frac{16n}{\pi}(-1)^{n}\pp{\frac{12n^2-19}{(4n^2-5^2)(4n^2-9^2)}}\\
                &= (-1)^{n}\frac{16n(12n^2-19)}{\pi(4n^2-25)(4n^2-81)}.
            \end{align*}
            Por lo tanto,
            \[ u(x,t)=\sum_{n\in\N}(-1)^{n}\frac{16n(12n^2-19)}{\pi(4n^2-25)(4n^2-81)} \sin\pp{\frac{n\pi x}{L}}e^{-\pp{\frac{\a n\pi}{L}}^2t}.\]
            \eej
        \end{enumerate}
        

        \item Resolver las siguientes PDEs relacionadas a la ecuación de onda por los métodos de separación de variables y con el método de D’Alembert,

        \begin{enumerate}
            \item \[
            \begin{cases}
                u_{tt} &= c^2u_{xx},\quad 0<x<L,~t>0,\\
                u_x(0,t) &= 0,\\
                u_x(\ell, t) &= 0,\\
                u(x,0) &= \phi(x),\\
                u_t(x,0) &= \psi(x).
            \end{cases}
            \]

            \item \[
            \begin{cases}
                u_{tt} &= a^2u_{xx},\quad 0\leq x\leq \pi,~t>0,\\
                u(0,t) &= 0,\quad t>0,\\
                u(\pi, t) &= 0,\quad t>0,\\
                u(x,0) &= x(\pi-x),\quad 0\leq x\leq\pi,\\
                u_t(x,0) &= \sin(7\pi x),\quad 0\leq x\leq\pi.
            \end{cases}
            \]
        \end{enumerate}

        \sol Primero resolveremos el caso general (primer inciso) por los dos métodos. Luego, para el segundo inciso usaremos el caso general

        Por el método de separación de variables supongamos que $u(x,t)=v(x)w(t)$ con $v, w\neq 0$. Entonces,
        \[ u_x(0,t)=u_x(L,t)=0\Ssi\begin{cases}
            v'(0)w(t)=0\\
            v'(L)w(t)=0
        \end{cases}\Ssi  v'(0)=v'(L)=0 . \]
        Luego,
        \[ u_{tt} = a^2 u_{xx} \Ssi vw''=a^2 v''w\Ssi\frac{w''}{w}=a^2\frac{v''}{v}=\l, \]
        donde $\l$ puede ser cualquier constante ya que las variables $x$ y $t$ son independientes entre ellas.

        Ahora buscamos la solución general de $z''=Az$, supongamos que la solución es de la forma $z(\xi)=e^{r\xi}$, tenemos que
        \[ z''-Az=0\Ssi r^2e^{r\xi}-Ae^{r\xi}=0\Ssi r=\pm\sqrt{A}. \]
        Entonces $v(x)=Ae^{\tfrac{\sqrt{\l}}{|a|}x}+Be^{-\tfrac{\sqrt{\l}}{|a|}x}$, si $\l=0$, entonces $v''=0$ con solución general $v(x)=Ax+B$ y $v'(x)=A$, para que $v'(0)=v'(L)=0$ se debe cumplir que $A=0$ por lo cual $v$ es un función constante, pero como $u(x,0)=\phi(x)$, entonces $v(x)=\phi(x)/w(0)$ con $w(0)\neq 0$ por lo cual $v$ no puede ser constante. Si $\l>0$, entonces
        \[ 
        v'(0)=v'(L)=0\Ssi \begin{cases}
            A\tfrac{\sqrt{\l}}{|a|}-\tfrac{\sqrt{\l}}{|a|} B=0\\
            A\tfrac{\sqrt{\l}}{|a|}e^{\tfrac{\sqrt{\l}}{|a|}L}-\tfrac{\sqrt{\l}}{|a|} Be^{-\tfrac{\sqrt{\l}}{|a|}L} =0
        \end{cases}
        \ssi
        \begin{cases}
            A=B\\
            Ae^{\tfrac{\sqrt{\l}}{|a|}L} =Ae^{-\tfrac{\sqrt{\l}}{|a|}L}
        \end{cases},
        \]
        lo cual ocurre sólo cuando $A=B=0$ que implica $v=0$ lo cual no queremos. Por otro lado, si $\l<0$, entonce
        \[ v(x) = Ae^{i\frac{\sqrt{-\l}}{|a|}} + Be^{-i\frac{\sqrt{-\l}}{|a|}} = A\cos\pp{\frac{\sqrt{-\l}}{|a|}x}+ B\sin\pp{\frac{\sqrt{-\l}}{|a|}x}. \]
        y
        \[ v'(x) =-A\frac{\sqrt{-\l}}{|a|}\sin\pp{\frac{\sqrt{-\l}}{|a|}x}+ B\frac{\sqrt{-\l}}{|a|}\cos\pp{\frac{\sqrt{-\l}}{|a|}x}. \]
        Luego
        \[ v'(0)=v'(L)=0\ssi \begin{cases}
             B\tfrac{\sqrt{-\l}}{|a|}=0\\
            -A\sin\pp{\tfrac{\sqrt{-\l}}{|a|}L}+ B\cos\pp{\tfrac{\sqrt{-\l}}{|a|}L}=0
        \end{cases}\ssi\begin{cases}
            B=0\\
            \sin\pp{\tfrac{\sqrt{-\l}}{|a|}L}=0
        \end{cases}. \]
        Como queremos que $v\neq 0$ entonces $A\neq 0$ y po lo tanto sólo consideramos los $\l_n$ tales que $\tfrac{\sqrt{-\l_n}}{|a|}L=n\pi$ con $n\in\N$, estos son $\l_n=-(n\pi a/L)^2$ y
        \[ v_n(x)= A_n\cos\pp{\frac{\sqrt{-\l_n}}{|a|}x}= A_n\cos\pp{\frac{n\pi x}{L}}.\]
        Por otro lado, ${w_n}''=\l_nw_n$ tiene como solución general
        \[ w_n(t)=B_ne^{i\tfrac{n\pi |a|}{L}t}+ C_n e^{-i\tfrac{n\pi |a|}{L}t}= B_n\cos\pp{\tfrac{n\pi |a|}{L}t}+ C_n\sin\pp{\tfrac{n\pi |a|}{L}t}. \]

        Por lo tanto, podemos escribir la solución como
        \[ u(x,t) = \sum_{n\in\N}v_n(x)w_n(t)=\sum_{n\in\N}\cos\pp{\tfrac{n\pi x}{L}}\cc{A_n\cos\pp{\tfrac{n\pi |a|}{L}t}+ B_n\sin\pp{\tfrac{n\pi |a|}{L}t}}\]
        y
        \[ u_t(x,t) = \sum_{n\in\N}\cos\pp{\tfrac{n\pi x}{L}}\cc{-A_n \tfrac{n\pi |a|}{L}\sin\pp{\tfrac{n\pi |a|}{L}t}+ B_n\tfrac{n\pi |a|}{L}\cos\pp{\tfrac{n\pi |a|}{L}t}}.\]
        Luego,
        \[ \begin{cases}
                u(x,0)=\phi(x)\\
                u_t(x,0)=\psi(x)
            \end{cases}
            \Ssi
            \begin{cases}
                \sum_{n\in\N}A_n\cos\pp{\tfrac{n\pi x}{L}}=\phi(x)\\
                \sum_{n\in\N}B_n\tfrac{n\pi |a|}{L}\cos\pp{\tfrac{n\pi x}{L}}=\psi(x)
            \end{cases}. \]
        Por la expansión de la serie de cosenos de Fourier se tiene que
        \[ A_n = \frac{2}{L}\int_0^{L}\phi(x)\cos\pp{\tfrac{n\pi x}{L}}dx\Text{y}B_n= \frac{2}{n\pi|a|}\int_0^{L}\psi(x)\cos\pp{\tfrac{n\pi x}{L}}dx.\]

        Ahora resolveremos la ecuación usando el método de D'Alembert. Sean
        \[ \xi = x-at\Text{y}\eta=x+at. \]
        Usando la regla de la cadena dos veces tenemos
        \begin{align*}
            u_x =\frac{\partial u}{\partial\xi}\frac{\partial\xi}{\partial x}&+\frac{\partial u}{\partial\eta}\frac{\partial\eta}{\partial x}=u_\xi+u_\eta,\\
            \Imp u_{xx} =u_{\xi\xi}+u_{\xi\eta}&+u_{\eta\eta}+u_{\eta\xi}=u_{\xi\xi}+2u_{\xi\eta}+u_{\eta\eta},
        \end{align*}
        ya que
        \[ \frac{\partial\xi}{\partial x}=\frac{\partial\eta}{\partial x}=1. \]
        Por otro lado, de manera analoga tenemos
        \begin{align*}
            u_t = \frac{\partial u}{\partial\xi}\frac{\partial\xi}{\partial t}&+\frac{\partial u}{\partial\eta}\frac{\partial\eta}{\partial t}=-au_\xi+au_\eta,\\
            \Imp u_{tt} = a(au_{\xi\xi}-au_{\xi\eta}&-au_{\xi\eta}+au_{\eta\eta}) = a^2(u_{\xi\xi}-2u_{\xi\eta}+u_{\eta\eta}) 
        \end{align*}
        ya que
        \[ \frac{\partial\xi}{\partial t}=-a\Text{y} \frac{\partial\eta}{\partial t}=a. \]
        Sustituyendo en la ecuación original tenemos que
        \[ a^2(u_{\xi\xi}-2u_{\xi\eta}+u_{\eta\eta}) = a^2(u_{\xi\xi}+2u_{\xi\eta}+u_{\eta\eta})\Imp u_{\xi\eta}=0. \]
        Luego, itegrando dos veces tenemos
        \[ u(x,t)=\iint u_{\xi\eta}~d\eta d\xi=\int f(\xi)d\xi=f(\xi)+g(\eta)=f(x-at)+g(x+at),\]
        donde $f$ y $g$ son funciones arbitrarias.

        Luuego tenemos que
        \begin{align*}
            u(x, 0)=\phi(x),~u_t(x,0)=\psi(x) &\Ssi f(x)+g(x)=\phi(x),~ -af'(x)+ag'(x)=\psi(x)\\
            &\Ssi f'(x)+g'(x)=\phi'(x),~ g'(x)-f'(x)=\tfrac{\psi(x)}{a}.
        \end{align*}
        Sumando y restando ambas expresiones tenemos que
        \[ g'(x)=\tfrac{a\phi'(x)+\psi(x)}{2a}\Text{y} f'(x)= \tfrac{a\phi'(x)-\psi(x)}{2a}. \]
        Integrando tenemos que
        \[f(x)=\frac{1}{2}\phi(x)-\frac{1}{2a}\int\psi(x)dx\Text{y}g(x)=\frac{1}{2}\phi(x)+\frac{1}{2a}\int\psi(x)dx\]
        Por lo tanto
        \[u(x,t)=\frac{1}{2}\cc{\phi(x-at)+\phi(x+at)}+\frac{1}{2a}\int_{x-ct}^{x+ct}\psi(z)dz.\]
        

        

        
        
        
        
        \eej


        \item Para cada uno de los problemas grafica las soluciones con un programa en Python que calcule los $n$ primeros
        términos de las series de Fourier soluciones. Definir los valores para los dominios $L$ y $T$, para $x$ y $t$ respectivamente. Las constantes de las ecuaciones pueden ser variables. Grafica las soluciones con $n=1,2,5,10, 20$.
    \end{enumerate}
\end{document}
